\section{Introducción a la Inducción Estructural}

La inducción es una técnica poderosa utilizada para demostrar propiedades de estructuras definidas recursivamente. Aunque la inducción matemática es útil para propiedades de los números naturales, la \textbf{inducción estructural} extiende esta idea a conjuntos definidos recursivamente.

A veces es difícil definir un objeto explícitamente, pero no lo es tanto si lo definimos en términos de objetos de su mismo tipo. A este proceso lo llamamos \textbf{recursión}. Mediante recursión, podemos definir secuencias, funciones y conjuntos, especificando un caso base y reglas que permiten construir nuevos elementos a partir de los ya existentes.


\subsection{Definiciones Recursivas}
\begin{definitionblock}{Definición Recursiva}
Un objeto puede definirse recursivamente especificando:
\begin{itemize}
    \item \textbf{Caso base}: Elementos iniciales del conjunto.
    \item \textbf{Regla recursiva}: Forma de construir nuevos elementos a partir de los ya definidos.
\end{itemize}
\end{definitionblock}
\subsubsection*{Ejemplo: Secuencia de Potencias de 2}
\begin{itemize}
    \item Caso base: $a_0 = 1$.
    \item Regla recursiva: $a_{n+1} = 2a_n$.
\end{itemize}
\noindent Aplicando esta definición recursiva, obtenemos los términos de la secuencia que generan las potencias de 2: $1, 2, 4, 8, 16, 32, \dots$.

La ventaja de definir una secuencia recursivamente es que no necesitamos conocer una fórmula explícita general desde el principio. En muchos casos, esto facilita la comprensión y el diseño de algoritmos que dependen de estructuras previas, como en el caso de la programación dinámica o la resolución de problemas combinatorios.


\subsection{Definición Recursiva de Funciones}
\subsubsection{Factorial}
Un ejemplo clásico de función definida recursivamente es el factorial. La función $f(n) = n!$ puede ser definida recursivamente como:
\[
    f(0) = 1, \quad f(n+1) = (n+1)f(n).
\]

\bigskip
\noindent \textbf{Ejercicio}: Defina recursivamente la función $S(n) = \sum_{k=0}^{n} a_k$.

\subsubsection{Números de Fibonacci}
Los números de Fibonacci están definidos por:
\begin{align*}
    f_0 &= 0, \quad f_1 = 1, \\
    f_n &= f_{n-1} + f_{n-2}, \quad \text{para } n \geq 2.
\end{align*}

Ejemplo: $0, 1, 1, 2, 3, 5, 8, 13, 21, \dots$

\subsubsection*{Ejercicio: Propiedad de Crecimiento de Fibonacci}
Demuestre que para todo $n \geq 3$, si $\alpha = \frac{1 + \sqrt{5}}{2}$, entonces $f_n > \alpha^{n-2}$.
\subsubsection*{Demostración por inducción:}
\noindent \textbf{Caso Base:} Debemos demostrar el caso base tanto para $n=3$ como para $n=4$ (porque no podemos demostrar que $f_4 > \alpha^2$ desde la hipótesis inductiva).

Esto no es difícil puesto que:
\begin{align*}
    f_3 &= 2 > \alpha = \frac{1 + \sqrt{5}}{2}, \\
    f_4 &= 3 > \alpha^2 = \frac{3 + \sqrt{5}}{2}.
\end{align*}

\noindent \textbf{Caso Inductivo:} Considere un entero $n + 1$ tal que $n \geq 4$. Por hipótesis inductiva, $f_j > \alpha^{j-2}$ para todo $j \in [3, n]$.

Dado que $f_{n+1} = f_n + f_{n-1}$, se tiene:
\[
    f_{n+1} > \alpha^{n-2} + \alpha^{n-3} = \alpha^{n-3} (1 + \alpha).
\]

Por tanto, para demostrar que $f_{n+1} > \alpha^{n-1}$ basta demostrar que:
\[
    (1 + \alpha) \geq \alpha^2.
\]

Esto es fácil, ya que:
\[
    (1 + \alpha) = \frac{3 + \sqrt{5}}{2} = \alpha^2.
\]

Así, la desigualdad se cumple y la demostración por inducción queda completada.

\bigskip
\noindent \textbf{Ejercicio}: Demuestre que para todo entero positivo $n$, 
 $$ f_0 f_1 + f_1 f_2 + \cdots + f_{2n-1} f_{2n} = f_{2n}^2. $$


\subsection{Conjuntos Definidos Recursivamente}
Los conjuntos también pueden definirse de manera recursiva, siguiendo una estructura similar a la de las funciones. Un conjunto recursivamente definido se construye especificando un conjunto base y una regla que permite generar nuevos elementos a partir de los que ya están en el conjunto.

La definición recursiva de conjuntos sigue la estructura:
\begin{itemize}
    \item \textbf{Caso base}: Elementos iniciales del conjunto.
    \item \textbf{Regla recursiva}: Forma de construir nuevos elementos.
\end{itemize}

\begin{exampleblock} {Definición de $\Sigma^*$} 
El conjunto de strings sobre un alfabeto $\Sigma$:
\begin{itemize}
    \item \textbf{Caso base}: La cadena vacía $\epsilon$ pertenece a $\Sigma^*$.
    \item \textbf{Regla recursiva}: Si $w \in \Sigma^*$ y $x \in \Sigma$, entonces $wx \in \Sigma^*$.
\end{itemize}
\end{exampleblock}
Este conjunto, denotado como $\Sigma^*$, representa todas las posibles cadenas (incluyendo la cadena vacía) que pueden formarse con los símbolos de $\Sigma$.

\subsubsection{Definiciones Recursivas de Operaciones sobre Strings}
Al igual que los conjuntos, las operaciones sobre strings pueden definirse de manera recursiva. Definir operaciones de esta manera nos permite formalizar algoritmos y probar propiedades sobre strings de manera estructurada.

Una de las operaciones más importantes sobre strings es la concatenación, que une dos strings en un solo. La concatenación de dos strings $w_1$ y $w_2$ se define recursivamente de la siguiente manera:
\begin{exampleblock}{Definición de concatenación}
\begin{itemize}
    \item \textbf{Regla base}: Si $w \in \Sigma^*$, entonces $w \cdot \epsilon = w$ (concatenar con la cadena vacía no cambia la cadena original).
    \item \textbf{Regla inductiva}: Si $w_1, w_2 \in \Sigma^*$ y $x \in \Sigma$, entonces $w_1 \cdot (w_2x) = (w_1 \cdot w_2) x$ (concatenar un carácter a una cadena es equivalente a concatenar primero el prefijo y luego agregar el carácter final).
\end{itemize}
\end{exampleblock}

Por ejemplo, si tenemos el alfabeto $\Sigma = \{a, b\}$ y las cadenas $w_1 = ab$, $w_2 = ba$, la concatenación $w_1 \cdot w_2$ sigue:
\begin{align*}
    ab \cdot ba &= (ab \cdot b)a \\
    &= ((ab \cdot \epsilon)b)a \\
    &= (ab)b)a \\
    &= abb a.
\end{align*}

Otra propiedad importante de los strings es su longitud. Podemos definir recursivamente la función $l(w)$, que calcula la longitud de un string $w$:
\begin{exampleblock}{Definición de longitud}
\begin{itemize}
    \item \textbf{Regla base}: $l(\epsilon) = 0$.
    \item \textbf{Regla inductiva}: $l(w x) = l(w) + 1$, para $x \in \Sigma$ (cada vez que agregamos un carácter a la cadena, su longitud aumenta en uno).
\end{itemize}
\end{exampleblock}

Por ejemplo, si $w = ab$:
\begin{align*}
    l(\epsilon) &= 0, \\
    l(a) &= l(\epsilon) + 1 = 1, \\
    l(ab) &= l(a) + 1 = 2.
\end{align*}

\section{Inducción Estructural sobre Estructuras Recursivas}
La \textbf{inducción estructural} es un método para demostrar propiedades de conjuntos definidos de manera recursiva. Dado que estos conjuntos están construidos a partir de reglas que generan nuevos elementos a partir de los anteriores, la inducción estructural nos permite razonar sobre ellos en una forma natural y efectiva.
La inducción estructural permite demostrar que una propiedad se cumple para todos los elementos del conjunto sin necesidad de verificarlos uno por uno.

El método de inducción estructural sigue estos pasos:
\begin{definitionblock}{Inducción Estructural}
\begin{itemize}
    \item \textbf{Caso base}:  Demostrar que la propiedad se cumple para cada caso base de
la definición recursiva.
    \item \textbf{Paso inductivo}: Se asume que la propiedad es cierta para los elementos ya definidos (hipótesis inductiva) y se prueba que debe ser cierta para los nuevos elementos generados por la regla recursiva.
\end{itemize}
\end{definitionblock}

\subsection{Ejemplo: Longitud de Cadenas}

\textbf{Ejercicio:} Demuestre que para todo par de strings $w_1, w_2 \in \Sigma^*$, se cumple:
\[
    l(w_1 w_2) = l(w_1) + l(w_2).
\]

\subsubsection*{Demostración}
\noindent Definimos la proposición $P(y)$ como:
\[
    l(x y) = l(x) + l(y) \quad \text{para cualquier } x \in \Sigma^*.
\]
Queremos probar que $P(y)$ es cierta para todo $y \in \Sigma^*$ usando inducción estructural.

\noindent \textbf{Caso base:} Si $y = \epsilon$ (cadena vacía), entonces:
\[
    l(x \epsilon) = l(x) = l(x) + 0 = l(x) + l(\epsilon).
\]
Esto muestra que la propiedad se cumple para el caso base.

\noindent \textbf{Paso inductivo:} Supongamos que la propiedad es cierta para una cadena $w$, es decir, que:
\[
    l(x w) = l(x) + l(w).
\]
Ahora, consideremos una cadena extendida con un símbolo adicional $a \in \Sigma$, es decir, $w a$. Queremos demostrar que:
\[
    l(x w a) = l(x) + l(w a).
\]
Por la definición recursiva de la longitud de un string, sabemos que:
\[
    l(w a) = l(w) + 1.
\]
Aplicando la hipótesis inductiva:
\begin{align*}
    l(x w a) &= l(x w) + 1 \\
    &= (l(x) + l(w)) + 1 \\
    &= l(x) + (l(w) + 1) \\
    &= l(x) + l(w a).
\end{align*}
Dado que hemos demostrado que la propiedad se mantiene en el caso base y en el paso inductivo, por el principio de inducción estructural, concluimos que la propiedad es cierta para toda cadena $y \in \Sigma^*$.

\subsection{Definición Recursiva de Árboles Binarios}
\begin{exampleblock}{Árbol Binario}
Un árbol binario completo se define como:
\begin{itemize}
    \item \textbf{Caso base:} Un árbol con un único nodo raíz $r$.
    \item \textbf{Regla recursiva:} Si $T_1$ y $T_2$ son árboles binarios completos, entonces $T = T_1 \cdot T_2$ es un árbol binario completo que consiste de una raíz r junto con arcos que unen a $r$ con la raíz del sub-árbol izquierdo $T_1$ y con la raíz del sub-árbol derecho $T_2$.
\end{itemize}
\end{exampleblock}

\subsubsection{Ejercicio: Número de Vértices en un Árbol Binario}
Defina recursivamente la función $n(T)$ que entrega el número de vértices de un árbol binario $T$.

\textbf{Solución:}
\begin{itemize}
    \item \textbf{Caso base}: Si $T$ está formado sólo por una raíz $r$, entonces $n(T) = 1$.
    \item \textbf{Regla inductiva}: Si $T_1$ y $T_2$ son árboles binarios completos, entonces el árbol binario completo $T = T_1 \cdot T_2$ tiene:
    \[
        n(T) = 1 + n(T_1) + n(T_2).
    \]
\end{itemize}

\subsubsection{Ejercicio: Altura de un Árbol Binario}
Defina la función $h(T)$ que entrega la distancia máxima entre la raíz de un árbol binario $T$ y una de sus hojas.

\textbf{Solución:}
\begin{itemize}
    \item \textbf{Caso base}: Si $T$ está formado sólo por una raíz $r$, entonces $h(T) = 0$.
    \item \textbf{Regla inductiva}: Si $T_1$ y $T_2$ son árboles binarios completos, entonces el árbol binario completo $T = T_1 \cdot T_2$ tiene:
    \[
        h(T) = 1 + \max(h(T_1), h(T_2)).
    \]
\end{itemize}

\subsubsection{Ejercicio: Relación entre $n(T)$ y $h(T)$}
Demuestre que para todo árbol binario completo $T$, se cumple que:
\[
    n(T) \leq 2^{h(T)+1} - 1.
\]

\textbf{Demostración por Inducción Estructural:}

\textbf{Caso base:} Si $T$ está formado sólo por una raíz $r$, entonces $n(T) = 1$ y $h(T) = 0$, por lo que:
\[
    n(T) = 1 \leq 2^{0+1} - 1 = 1.
\]

\textbf{Paso inductivo:} Supongamos que la propiedad se cumple para los subárboles $T_1$ y $T_2$, es decir,
\[
    n(T_1) \leq 2^{h(T_1)+1} - 1, \quad n(T_2) \leq 2^{h(T_2)+1} - 1.
\]

% Por las fórmulas recursivas de $n(T)$ y $h(T)$, tenemos:
% \[
%     n(T) = 1 + n(T_1) + n(T_2),
% \]
% \[
%     h(T) = 1 + \max(h(T_1), h(T_2)).
% \]

Utilizando la hipótesis inductiva:
\begin{align*}
n(T) &= 1 + n(T_1) + n(T_2) \quad  \\
&\leq 1 + (2^{h(T_1)+1 }- 1) + (2^{h(T_2)+1} - 1) \quad \text{por H.I.} \\
&\leq 2 \cdot \max(2^{h(T_1)+1} , 2^{h(T_2)+1}) - 1 \quad \text{suma es como máx 2 veces el mayor} \\
&= 2 \cdot 2^{max(h(T_1),h(T_2))+1} - 1 \quad \text{ } max(2^x,2^y) = 2^{max(x,y)} \\
&= 2 \cdot 2^{h(T)+1} - 1 \quad \text{por def. de } h(T) \\
&= 2^{h(T)+1} - 1.
\end{align*}


Por lo tanto, la propiedad se cumple para todo árbol binario completo $T$, completando la demostración.


% \section{Conclusión}
% La inducción estructural es una técnica esencial en matemáticas discretas y ciencias de la computación. Permite demostrar propiedades de estructuras definidas recursivamente, como funciones, secuencias, conjuntos y árboles binarios.


\subsection{Ejercicios Propuestos:}
\begin{itemize}
    % \item Demostrar que la concatenación de dos cadenas en $\Sigma^*$ es asociativa, es decir, para todo $x, y, z \in \Sigma^*$:
    % \[
    %     (x \cdot y) \cdot z = x \cdot (y \cdot z).
    % \]
     \item Definir recursivamente la función $rev(w)$ que invierte un string en $\Sigma^*$, y demostrar por inducción estructural que $rev(xy) = rev(y)rev(x)$ para todo $x, y \in \Sigma^*$.
    \item Sea $A$ el conjunto definido recursivamente por: \textbf{(i)} $\epsilon \in A$; \textbf{(ii)} si $w \in A$ entonces $awa \in A$, para $a \in \Sigma$. Demostrar que toda cadena en $A$ tiene longitud par.

    \item Definir recursivamente la función que cuenta el número de hojas en un árbol binario completo.
    \item Probar por inducción estructural que el número de hojas en un árbol binario completo de altura $h$ es exactamente $2^h$.
    \item Demostrar que en un árbol binario completo con $n$ nodos internos, hay exactamente $n+1$ hojas.
    \item Probar que la altura de un árbol binario completo con $n$ nodos es al menos $\log_2(n+1) - 1$.
\end{itemize}


\section{Ejercicios Resueltos}

\subsection{Ejercicio: Propiedad de la reversa de strings}

Definir recursivamente la función $rev(w)$ que invierte un string $w \in \Sigma^*$. Luego, demostrar por inducción estructural que:
\[
    \text{rev}(xy) = \text{rev}(y)\text{rev}(x)
\]
para todo $x, y \in \Sigma^*$.

\textbf{Definición recursiva de $\text{rev}(w)$:}
\begin{itemize}
    \item Caso base: $\text{rev}(\epsilon) = \epsilon$
    \item Paso recursivo: Si $w \in \Sigma^*$ y $a \in \Sigma$, entonces $\text{rev}(wa) = a \cdot \text{rev}(w)$
\end{itemize}

\textbf{Demostración:} Fijamos un string $x \in \Sigma^*$. Demostraremos por inducción estructural sobre $y$ que:
\[
    \text{rev}(xy) = \text{rev}(y) \cdot \text{rev}(x)
\]

\textbf{Caso base:} $y = \epsilon$
\[
    \text{rev}(x\epsilon) = \text{rev}(x) = \epsilon \cdot \text{rev}(x) = \text{rev}(\epsilon) \cdot \text{rev}(x)
\]

\textbf{Paso inductivo:} Supongamos que $\text{rev}(xw) = \text{rev}(w)\text{rev}(x)$ para una cadena $w \in \Sigma^*$. Queremos demostrar que también se cumple para $wa$, donde $a \in \Sigma$.
\begin{align*}
\text{rev}(xwa) &= a \cdot \text{rev}(xw) \\
                &= a \cdot (\text{rev}(w)\text{rev}(x)) \\
                &= (a \cdot \text{rev}(w)) \cdot \text{rev}(x) \\
                &= \text{rev}(wa) \cdot \text{rev}(x)
\end{align*}
Por lo tanto, se cumple la propiedad para todo $y \in \Sigma^*$.

\subsection{Ejercicio: Longitud par en un conjunto definido recursivamente}

Sea el conjunto $A$ definido por:
\begin{itemize}
    \item Caso base: $\epsilon \in A$
    \item Paso recursivo: Si $w \in A$ y $a \in \Sigma$, entonces $awa \in A$
\end{itemize}

Demostrar que toda cadena en $A$ tiene longitud par.

\textbf{Demostración por inducción estructural:}

\textbf{Caso base:} $\epsilon \in A$, y su longitud es $0$, que es par.

\textbf{Paso inductivo:} Supongamos que $|w|$ es par. Entonces $|awa| = |w| + 2$. Como $|w|$ es par, $|w| = 2k$ para algún $k$, por lo tanto:
\[
    |awa| = 2k + 2 = 2(k + 1)
\]
que es par.

Por lo tanto, toda cadena en $A$ tiene longitud par.